\section{Model}
\label{sec:model}

\subsection{Tree-Based Data Trading Game}
\label{subsec:game}

We consider a data trading game on a rooted tree $\calT$.
The node set $\calI$ is the set of all players, with each directed edge pointing from a parent to a child. Let $r \in \calI$ denote the root, and $\calL \subseteq \calI$ denote the set of leaf nodes. We write $N \coloneqq |\calI|$ for the total number of players. We define the following for players:
\begin{itemize}[left=0em]
\item $\Parent(i)$, for $i \in \calI \setminus \{r\}$: the unique parent of $i$ in $\calT$.
\item $\Children(i)$, for $i \in \calI$: the set of children of $i$. For leaf nodes $i \in \calL$, $\Children(i) = \varnothing$.
\item $\Path(i)$, for $i \in \calI$: the set of strict ancestors of $i$, \ie, the nodes on the path from $r$ to $i$ excluding $i$. We have $\Path(r) = \varnothing$.
\item $\Subtree(i)$, for $i \in \calI$: all nodes in the subtree rooted at $i$, including $i$ itself.
\item $\Desc(i) \coloneqq \Subtree(i) \setminus \{i\}$, for $i \in \calI$: the set of strict descendants of $i$.
\end{itemize}
\cref{fig:tree-notation} illustrates these notations on a small instance.

\begin{figure}[t]
\centering
\resizebox{0.66\textwidth}{!}{%
\begin{tikzpicture}[
  >={Latex[length=1.8mm]},
  vtx/.style={circle, draw, line width=0.5pt, minimum size=7.5mm, inner sep=0pt, font=\small},
  rt/.style={vtx, fill=croot, text=white},
  lf/.style={vtx, fill=cleaf, text=white},
  rf/.style={vtx, fill=cref},
  it/.style={vtx, fill=white},
  ed/.style={->, line width=0.6pt, black!70},
]
% tree
\node[rt] (r)   at (0,0)        {$r$};
\node[lf] (l)   at (-1.5,-1.35) {$l$};
\node[it] (u)   at (1.5,-1.35)  {$u$};
\node[rf] (ul)  at (0.55,-2.75) {$ul$};
\node[lf] (uu)  at (2.65,-2.75) {$uu$};
\node[lf] (ull) at (-0.35,-4.1) {$ull$};
\node[lf] (ulu) at (1.45,-4.1)  {$ulu$};
\draw[ed] (r)--(l);  \draw[ed] (r)--(u);
\draw[ed] (u)--(ul); \draw[ed] (u)--(uu);
\draw[ed] (ul)--(ull); \draw[ed] (ul)--(ulu);
% subtree box
\begin{scope}[on background layer]
\node[draw=cref, dashed, line width=1pt, rounded corners=3pt,
      fit=(ul)(ull)(ulu), inner sep=7pt] (box) {};
\end{scope}
\node[text=cref, font=\small, anchor=south west] at ([xshift=-3pt,yshift=1.5pt]box.north west) {$\Subtree(ul)$};
% legend
\node[rt] (lr) at (5.0,0.05) {$r$};
\node[right=1.2mm of lr, font=\small] {: root};
\node[lf] (ll) at (5.0,-1.05) {$l$};
\node[right=1.2mm of ll, font=\small] {: leaves};
% info brace
\node[anchor=west, font=\small] (info) at (4.55,-3.05)
  {$\left\{\begin{array}{l}\Parent(ul)=u\\[5pt]\Children(ul)=\{ulu,ull\}\\[5pt]\Path(ul)=\{r,u\}\end{array}\right.$};
\draw[->, dashed, line width=0.7pt, black!70] (ul.south east) to[out=-38,in=182] (info.west);
\end{tikzpicture}%
}
\caption{An illustration of the tree notation on a $7$-player instance. The root $r$ is shown in blue and the leaves $\calL = \{uu, l, ulu, ull\}$ in green. Taking node $ul$ as the reference: its parent is $\Parent(ul) = u$, its children are $\Children(ul) = \{ulu, ull\}$, and its strict ancestors are $\Path(ul) = \{r, u\}$. The red dashed box marks $\Subtree(ul) = \{ul, ulu, ull\}$, so that $\Desc(ul) = \{ulu, ull\}$. Best viewed in color.}
\label{fig:tree-notation}
\end{figure}

Replicable data flow downward along the edges of $\calT$, originating at the root $r$. Once a non-root player $i$ receives data from its parent, it refines the data into a sellable product that it may, in turn, offer to each of its children. Player $i$'s \emph{valuation} is a single scalar $v_i \in \calV_i \subseteq \bbR$: the net worth of the data product to $i$, namely the benefit it derives minus the processing cost, which may be negative. This quantity is observed only by player $i$.
We write $\bmv =\{v_i\}_{i \in \calI}$ as valuation profile, $\calV \coloneqq \prod_{i \in \calI} \calV_i$ for the space of valuation profiles, $\bmv_{-i}$ for the valuation profile other than player $i$'s, and $\calV_{-i} \coloneqq \prod_{k \neq i} \calV_k$ for its domain. We assume each $\calV_i$ is a closed bounded interval.
We let the valuation profile evolve as a Markov process down the tree.
%  to the tree, and because a refined product inherits its quality from the input it was built upon, we let valuations evolve as a Markov process down the tree.

\begin{assumption}[Markovian on Trees]
\label{asp:markov}
The root valuation is drawn as $v_r \sim f_r(\cdot)$. For every non-root player $i \in \calI \setminus \{r\}$, the valuation $v_i$ depends only on $v_{\Parent(i)}$:
\[
v_i \sim f_i(\cdot \mid v_{\Parent(i)}).
\]
The valuation profile $\bmv = \{v_i\}_{i \in \calI}$ has the joint distribution $f(\bmv) = f_r(v_r) \prod_{i\ne r} f_i(v_i|v_{\Parent(i)})$.
We denote the joint distribution of $\bmv$ by $\calF$, and the conditional CDF by $F_i$. Both $\calF$, $F_i$s and $f_i$s are common knowledge.
\end{assumption}

This assumption fits the tree structure: a player $i$'s product's worth is governed by the immediate input (player $\Parent(i)$'s product) from which it is refined, rather than by the raw material that forms the immediate input. It also subsumes the classical independent-private-values setting~\citep{myerson1981optimal} as a special case.


\subsection{Player Action}
\label{subsec:action}

Each edge $(i,j)$ with $j \in \Children(i)$ is a potential trade. The seller $i$ first commits to a posted price $p_j \in \bbR_+$ for child $j$; seeing this price, player $j$ responds with a binary decision $x_j \in \{0,1\}$, describing her willingness to buy ($x_j = 1$) or not ($x_j = 0$).

A trade $i$ actually occurs if and only if trade $\Parent(i)$ occurs and player $i$ is willing to buy the data from player $\Parent(i)$. We define the indicator of whether trade $i$ occurs:
\begin{equation}
\label{eq:yi}
y_i \;\coloneqq\; \prod_{j \in \Path(i) \cup \{i\}} x_j,
\end{equation}
where we set $x_r \coloneqq 1$ and $p_r \coloneqq 0$ by convention (the root player initially possesses the data).
If $y_i = 1$, trade $i$ occurs, player $i$ obtains the data, realizes valuation $v_i$, and pays $p_i$; otherwise $y_i = 0$, player $i$ gains nothing and pays nothing. For $m \in \Desc(j)$, we define the \emph{conditional trade indicator}:
\begin{equation}
\label{eq:y-conditional}
y_{m \mid j} \;\coloneqq\; \prod_{\substack{\ell \in \Path(m) \cup \{m\} \\ \ell \notin \Path(j) \cup \{j\}}} x_\ell,
\end{equation}
which is the product of trade decisions along the path from $j$ to $m$, excluding $j$. This captures whether trade $m$ occurs conditional on trade $j$ occurs.


\subsection{Player Strategy}
\label{subsec:strategy}

At each stage, a player observes the trading history along their path from the root. For player $i$, the observed history is:
\begin{equation}
\label{eq:hi}
h_i = \big\{(j, p_j, x_j)\big\}_{j \in \Path(i) \setminus \{r\}} \cup \{(i, p_i)\},
\end{equation}
which includes all past trade outcomes on the path from $r$ to $i$ and the price $p_i$ currently offered to $i$. We set $h_r = \varnothing$. Let $\calH_i$ denote the set of all possible histories for player $i$.

A strategy for player $i$ consists of:
\begin{itemize}[left=0em]
\item \textbf{Pricing strategy} (for $i \notin \calL$): $\tp_i: \calH_i \times \Children(i) \times \calV_i \to \bbR_+$, where $\tp_i(h_i, j, v_i)$ is the price player $i$ offers to child $j$.
\item \textbf{Buying strategy} (for $i \neq r$): $\tx_i: \calH_i \times \calV_i \to \{0, 1\}$, where $\tx_i(h_i, v_i)$ is the player $i$'s willingness to buy given public history and her private valuation $v_i$.
\end{itemize}
We call $(\tbmp, \tbmx) = \big(\{\tp_i\}_{i \notin \calL},\; \{\tx_i\}_{i \neq r}\big)$ a strategy profile.


\subsection{Game Dynamic}
\label{subsec:dynamic}

\begin{definition}[Game Dynamics on Trees]
\label{def:game-evolve}
Given a valuation profile $\bmv$ and a strategy profile $(\tbmp, \tbmx)$, the game evolves recursively from root to leaves:\\
1. Initialize $h_r = \varnothing$.\\
2. For each non-leaf player $i$ and each child $j \in \Children(i)$, the seller posts price $p_j = \tp_i(h_i, j, v_i)$.\\
3. Each non-root player $i$ decides $x_i = \tx_i(h_i, v_i)$.\\
4. The game concludes with the terminal history $h = \{(i, p_i, x_i)\}_{i \in \calI \setminus \{r\}}$.
\end{definition}

We denote this mapping from valuation profile $\bmv$, strategy profile $(\tbmp, \tbmx)$ to the terminal history $h$ as $h(\tbmp, \tbmx;\bmv)$.


\subsection{Profit Reallocation Mechanism}
\label{subsec:mechanism}

A \emph{profit reallocation mechanism} (PRM) specifies how the payment from each trade is distributed among upstream players.

\begin{definition}[Profit Reallocation Mechanism on Trees]
\label{def:prm}
A PRM is a function $\pi: (\calI \setminus \{r\}) \times (\calI \setminus \calL) \to \bbR_+$, where $\pi(i,j)$ is the proportion of the payment from trade $i$ allocated to player $j$. That is, when trade $i$ occurs at price $p_i$, player $j$ receives the monetary transfer of $\pi(i,j) \cdot p_i$. The function $\pi$ must satisfy following constraints to make it economically meaningful:
\begin{enumerate}[left=0em]
\item \textbf{Path Only:} $\pi(i,j) > 0$ only if $j \in \Path(i)$. Only ancestors of buyer $i$ are eligible for profit reallocation.
\item \textbf{Budget Feasibility:} $\sum_{j \in \Path(i)} \pi(i,j) \le 1$ for all $i \in \calI \setminus \{r\}$. The revenue of ancestors cannot exceed the payment from the buyer in any trade.
\end{enumerate}
\end{definition}

The PRM-aware model is specified by $\Model = (\calT, \pi, \calF)$.
A distinguished PRM is the \emph{baseline mechanism}, which does not perform profit reallocation:

\begin{definition}[Baseline Mechanism]
\label{def:baseline}
The \emph{baseline mechanism} $\pi^0$ allocates the entire payment to the direct seller: $\pi^0(i, \Parent(i)) = 1$ and $\pi^0(i, j) = 0$ for $j \neq \Parent(i)$.
\end{definition}


\subsection{Utility and Solution Concept}
\label{subsec:solution}

Given a PRM $\pi$ and terminal history $h$, the realized utility of player $i$ is:
\begin{equation}
\label{eq:utility-raw}
U_i(h; v_i) = y_i \cdot (v_i - p_i) + \sum_{j \in \Desc(i)} y_j \cdot \pi(j,i) \cdot p_j.
\end{equation}
The first term is player $i$'s own surplus from trade $i$; the second term aggregates the profit reallocation routed to $i$ from descendant trades, which by the Path Only condition (\cref{def:prm}) arrive only from $\Desc(i)$. We also denote social welfare as the aggregate utility $\SW(h;\bmv) = \sum_{i \in \calI} U_i(h;v_i)$.

\begin{figure*}[t]
\centering
\resizebox{0.9\textwidth}{!}{%
\begin{tikzpicture}[
  >={Latex[length=2mm]},
  rootN/.style={circle, draw, minimum size=8mm, inner sep=0pt, fill=croot,   text=white, font=\small},
  playerN/.style={circle, draw, minimum size=8mm, inner sep=0pt, fill=cplayer, text=white, font=\small},
  leafN/.style={circle, draw, minimum size=8mm, inner sep=0pt, fill=cleaf,   text=white, font=\small},
  db/.style={cylinder, shape border rotate=90, draw, aspect=0.28,
             minimum width=6.5mm, minimum height=5.5mm, fill=cdb, font=\tiny, inner sep=1pt},
  dyn/.style={font=\footnotesize, text=cdyn, align=left, inner sep=1pt},
  val/.style={font=\footnotesize, inner sep=1pt},
  ubox/.style={draw, rounded corners=2pt, font=\footnotesize, inner sep=3.5pt},
  ed/.style={->, line width=0.6pt, black!70},
  reOne/.style={->, line width=1pt, red!80!black},
  reHalf/.style={->, line width=1pt, creall},
  dcon/.style={-{Latex[length=1.6mm]}, dotted, line width=0.7pt, black!55},
]
% nodes
\node[rootN]   (r)  at (0,0)        {$r$};
\node[playerN] (l)  at (-3.4,-2.1)  {$l$};
\node[playerN] (u)  at (3.4,-2.1)   {$u$};
\node[leafN]   (ll) at (-5.1,-4.3)  {$ll$};
\node[leafN]   (lu) at (-2.1,-4.3)  {$lu$};
\node[leafN]   (ul) at (2.1,-4.3)   {$ul$};
\node[leafN]   (uu) at (5.1,-4.3)   {$uu$};
% data cylinders
\node[db, left=1.5mm of r]  (dbr) {data};
\node[db, left=1.5mm of u]  (dbu) {data};
\node[db, left=1.5mm of uu] {data};
% tree edges
\draw[ed] (r)--(l); \draw[ed] (r)--(u);
\draw[ed] (l)--(ll); \draw[ed] (l)--(lu);
\draw[ed] (u)--(ul); \draw[ed] (u)--(uu);
% valuations (internal nodes to the outer side to clear the child edges)
\node[val, below=0.5mm of r]  {$v_r=0$};
\node[val, anchor=east] at ([xshift=-1mm]l.west)   {$v_l=2$};
\node[val, anchor=east] at ([xshift=-1mm]dbu.west) {$v_u=-1$};
\node[val, below=0.5mm of ll] {$v_{ll}=3$};
\node[val, below=0.5mm of lu] {$v_{lu}=4$};
\node[val, below=0.5mm of ul] {$v_{ul}=5$};
\node[val, below=0.5mm of uu] {$v_{uu}=6$};
% game-dynamic labels (price offered / decision)
\node[dyn, right=1mm of l]  {$p_l=2$\\$x_l=0$};
\node[dyn, right=1mm of u]  {$p_u=1$\\$x_u=1$};
\node[dyn, right=1mm of ll] {$p_{ll}=6$\\$x_{ll}=0$};
\node[dyn, right=1mm of lu] {$p_{lu}=5$\\$x_{lu}=1$};
\node[dyn, right=1mm of ul] {$p_{ul}=4$\\$x_{ul}=0$};
\node[dyn, right=1mm of uu] {$p_{uu}=5$\\$x_{uu}=1$};
% profit-reallocation arrows (both orange arrows share the same start point on uu)
\coordinate (os) at (uu.north);
\draw[reOne]  (u.north west) to[out=150,in=-32]
      node[pos=0.5, right=0.6mm, font=\scriptsize, text=red!80!black] {$1$} (r.east);
\draw[reHalf] (os) to[out=122,in=-42]
      node[pos=0.42, left=0.4mm, font=\scriptsize, text=creall] {$0.5$} (u.south east);
\draw[reHalf] (os) to[out=76,in=-14]
      node[pos=0.10, right=0.4mm, font=\scriptsize, text=creall] {$0.5$} (r.east);
% utility boxes
\node[ubox, anchor=west] (Ur)  at (7.6,0.1)   {$U_r={\color{black}0}+{\color{red!80!black}1}\!\cdot\!{\color{cdyn}1}+{\color{creall}0.5}\!\cdot\!{\color{cdyn}5}=3.5$};
\node[ubox, anchor=west] (Uu)  at (7.6,-2.4)  {$U_u={\color{black}-1}-{\color{cdyn}1}+{\color{creall}0.5}\!\cdot\!{\color{cdyn}5}=0.5$};
\node[ubox, anchor=west] (Uuu) at (7.6,-4.3)  {$U_{uu}={\color{black}6}-{\color{cdyn}5}=1$};
\draw[dcon] (Ur.west)  to[bend right=8] (r.east);
\draw[dcon] (Uu.west)  -- (u.east);
\draw[dcon] (Uuu.west) -- (uu.east);
% brace + boxes for zero utilities
\draw[decorate, decoration={brace, amplitude=5pt, mirror}, black!70]
      ([yshift=-6mm]ll.south) -- ([yshift=-6mm]lu.south);
\node[ubox, below=13mm of $(ll)!0.5!(lu)$] {$U_l=U_{ll}=U_{lu}=0$};
\node[ubox, below=9mm of ul] {$U_{ul}=0$};
% legend
\begin{scope}[shift={(-5.3,-6.65)}]
  \node[rootN, scale=0.8] (lgp) at (0,0) {$r$};
  \node[right=1mm of lgp, font=\footnotesize] (lgpt) {: player};
  \node[right=6mm of lgpt, font=\footnotesize] (lgv) {$v_r$: private valuation};
  \node[right=6mm of lgv, font=\footnotesize, text=cdyn] (lgd) {$p_u,x_u$};
  \node[right=1mm of lgd, font=\footnotesize] (lgdt) {: game dynamic};
  \coordinate (a0) at ([xshift=7mm]lgdt.east);
  \draw[reHalf] (a0) -- ($(a0)+(0.8,0)$) node[midway, above=0.3mm, font=\scriptsize, text=creall] {$0.5$};
  \coordinate (a1) at ($(a0)+(1.25,0)$);
  \draw[reOne] (a1) -- ($(a1)+(0.8,0)$) node[midway, above=0.3mm, font=\scriptsize, text=red!80!black] {$1$};
  \node[anchor=west, font=\footnotesize] at ($(a1)+(0.95,0)$) {: profit reallocation};
\end{scope}
\end{tikzpicture}%
}
\caption{A running example of the tree-based data trading game. Best viewed in color.}
\label{fig:example}
\end{figure*}

\begin{example}
\label{ex:running}
\cref{fig:example} shows an instance with $N = 7$ players; the valuations, the PRM $\pi$, and all actions are given there, and the realized utilities follow from \cref{eq:utility-raw}. Two effects are worth highlighting. First, player $u$ has a \emph{negative} value $v_u = -1$, so its direct trade loses $v_u - p_u = -2$; yet the reallocated profit $\pi(uu,u)\,p_{uu} = 2.5$ from the successful downstream trade $uu$ leaves it with $U_u = 0.5$---so it's sometimes rational to buy the data even if the valuation is negative. Second, trade $lu$ fails despite $x_{lu}=1$, because the upstream trade $l$ fails ($x_l=0$), so $l, lu, ll$ all earn zero.
\end{example}

Under the strategy profile $(\tbmp, \tbmx)$, the expected utility of player $i$ is:
\begin{equation}
\label{eq:expected-utility}
u_i(\tbmp, \tbmx) \;\coloneqq\; \bbE_{\bmv \sim \calF}\!\big[U_i(h(\tbmp, \tbmx; \bmv); v_i)\big],
\end{equation}
where $h(\tbmp, \tbmx; \bmv)$ is the terminal history generated by $(\tbmp, \tbmx)$ and $\bmv$ according to game dynamic (\cref{def:game-evolve}).

How would strategic players behave under a given PRM? As the game is sequential and players hold private information, we adopt \emph{perfect Bayesian equilibrium} (PBE)~\citep{fudenberg1991game,MWG}---the standard refinement of Nash equilibrium for such games---as our solution concept: each player must act optimally at \emph{every} information set she may reach, under a belief about the others' valuations that is updated by Bayes' rule.

Player $i$'s \emph{information set} is the pair $(h_i, v_i)$ of her observed history and private valuation. A \emph{belief system} $\mu = \{\mu_i\}_{i \in \calI}$ specifies, at each information set, a distribution $\mu_i(\cdot \mid h_i, v_i) \in \Delta(\calV_{-i})$ over the other players' valuations $\bmv_{-i}$, where $\Delta(\calV_{-i})$ is the set of probability distributions on $\calV_{-i}$; this is player $i$'s posterior about $\bmv_{-i}$ upon observing $(h_i, v_i)$. Given a strategy profile $(\tbmp, \tbmx)$ and the belief $\mu_i$, player $i$'s conditional expected utility at $(h_i, v_i)$ is
\begin{equation}
\label{eq:cond-utility}
\bbU_i\big((\tbmp, \tbmx), \mu_i \mid h_i, v_i\big) \;\coloneqq\; \bbE_{\bmv_{-i} \sim \mu_i(\cdot \mid h_i, v_i)}\!\big[\, U_i\big(h(\tbmp, \tbmx;\, (v_i, \bmv_{-i})); v_i\big)\big],
\end{equation}
where $U_i$ is the realized utility~\eqref{eq:utility-raw} and $h(\cdot\,; \bmv)$ the terminal history generated by the game dynamics (\cref{def:game-evolve}).

\begin{definition}[Perfect Bayesian Equilibrium~\citep{fudenberg1991game,MWG}]
\label{def:pbe}
An assessment $\big((\tbmp^*, \tbmx^*), \tbmmu\big)$ of a strategy profile and a belief system is a \emph{perfect Bayesian equilibrium} if both of the following hold.
\begin{enumerate}[label=(\alph*),left=0em]
\item \textup{(Sequential rationality)} For every player $i \in \calI$, every information set $(h_i, v_i)$, and every deviation $(\tp_i, \tx_i)$ in player $i$'s own strategy components ($\tp_i$ only for the root, $\tx_i$ only for a leaf, both otherwise),
\[
\bbU_i\big((\tbmp^*, \tbmx^*), \tmu_i \mid h_i, v_i\big) \;\ge\; \bbU_i\big((\tp_i, \tbmp^*_{-i}, \tx_i, \tbmx^*_{-i}), \tmu_i \mid h_i, v_i\big).
\]
\item \textup{(Bayes-consistency)} At every information set $(h_i, v_i)$ reached with positive probability under $(\tbmp^*, \tbmx^*)$, the belief $\tmu_i(\cdot \mid h_i, v_i)$ equals the posterior distribution of $\bmv_{-i}$ induced by the prior $\calF$ and the profile $(\tbmp^*, \tbmx^*)$ given $(h_i, v_i)$.
\end{enumerate}
\end{definition}

% The Markovian structure makes them especially clean: by \cref{fact:simplify} below, the only component of $\mu_i$ that affects player $i$'s utility---the law of the descendant valuations $\bmv_{\Desc(i)}$ entering $U_i$ in~\eqref{eq:utility-raw}---equals $\calF(\cdot \mid v_i)$ at \emph{every} information set, on or off path. The payoff-relevant beliefs are thus pinned down everywhere, leaving none of the off-path latitude that ordinarily makes weak PBE permissive.\footnote{The Kreps--Wilson sequential-equilibrium refinement~\citep{kreps1982sequential} targets exactly this latitude, but is formulated for finite games; with continuous prices and valuations we do not invoke it, though the uniqueness just noted is the robustness it enforces.} To exhibit a PBE it therefore suffices to find a profile under which, at each $(h_i, v_i)$, player $i$'s prescribed action maximizes $\bbU_i$ with $\bmv_{\Desc(i)} \sim \calF(\cdot \mid v_i)$---exactly what the backward induction of \cref{sec:algo} establishes.
